\documentclass{article}
\usepackage[utf8]{inputenc}
\usepackage[T2A]{fontenc}
\usepackage[russian]{babel}
\usepackage{listings}
\usepackage{hyperref}

\title{da\_tune_scan: Программа для сканирования динамической апертуры на сетке бетатронных частот}
\author{}
\date{}

\begin{document}

\maketitle

\section{Введение}

Программа \texttt{da\_tune_scan} представляет собой комбинацию функциональности двух существующих программ BSim:
\begin{itemize}
    \item \texttt{tune\_scan} -- сканирование по сетке бетатронных частот с подстройкой lattice
    \item \texttt{dynamic\_aperture} -- расчет динамической апертуры методом трекинга
\end{itemize}

Основная цель программы -- исследовать зависимость динамической апертуры ускорителя от рабочих частот.

\section{Алгоритм работы}

\begin{enumerate}
    \item Чтение параметров из входного файла (по умолчанию \texttt{da\_tune_scan.init})
    \item Инициализация lattice и расчет равновесной орбиты
    \item Цикл по сетке бетатронных частот $(Q_a, Q_b, Q_z/pz)$:
    \begin{enumerate}
        \item Подстройка lattice для достижения заданных частот (используя \texttt{set\_tune\_3d})
        \item Пересчет равновесной орбиты и функций Твисса
        \item Расчет динамической апертуры в данной точке:
        \begin{itemize}
            \item Сканирование по углам от \texttt{min\_angle} до \texttt{max\_angle}
            \item Для каждого угла -- поиск радиуса потери устойчивости методом дихотомии
            \item Трекинг частицы на \texttt{n\_turn} оборотов
        \end{itemize}
        \item Вычисление площади DA и сохранение результатов
    \end{enumerate}
    \item Запись результатов в выходной файл
\end{enumerate}

\section{Структура данных}

\subsection{Тип da\_ts\_params\_struct}

Основные параметры:
\begin{lstlisting}[language=Fortran]
type da_ts_params_struct
  character(100) :: lat_file           ! Файл lattice
  character(100) :: dat_out_file       ! Файл результатов
  character(40)  :: quad_mask          ! Маска квадруполей
  character(40)  :: group_knobs(2)     ! Группы ручек настройки
  
  ! Сетка частот
  real(rp) :: Q_a0, Q_a1, dQ_a         ! Горизонтальная плоскость
  real(rp) :: Q_b0, Q_b1, dQ_b         ! Вертикальная плоскость
  real(rp) :: Q_z0, Q_z1, dQ_z         ! Продольная частота (RF on)
  real(rp) :: pz0, pz1, dpz            ! Энергетический разброс (RF off)
  
  logical :: rf_on                     ! Режим RF
  
  ! Параметры DA
  real(rp) :: min_angle, max_angle     ! Диапазон углов (радианы)
  integer :: n_angle                   ! Количество угловых точек
  integer :: n_turn                    ! Количество оборотов
  real(rp) :: x_init, y_init           ! Начальные амплитуды
  real(rp) :: abs_accuracy, rel_accuracy ! Точности
end type
\end{lstlisting}

\subsection{Тип da\_ts\_data\_struct}

Результаты расчета в каждой точке:
\begin{lstlisting}[language=Fortran]
type da_ts_data_struct
  integer :: ix_q(3)              ! Индексы точки сетки
  real(rp) :: tune(3)             ! Значения частот
  real(rp) :: da_area             ! Площадь динамической апертуры
  real(rp) :: da_x_max            ! Максимальная апертура по X
  real(rp) :: da_y_max            ! Максимальная апертура по Y
  integer :: n_angle_survived     ! Число выживших углов
  logical :: calc_successful      ! Флаг успеха
end type
\end{lstlisting}

\section{Расчет площади динамической апертуры}

Площадь вычисляется методом полярной интеграции:

\[
A = \frac{1}{2} \sum_{i=1}^{N-1} r_i^2 \Delta\theta
\]

где $r_i^2 = x_i^2 + y_i^2$ -- квадрат радиуса в i-м направлении, $\Delta\theta$ -- шаг по углу.

\section{Пример использования}

\subsection{Входной файл}

\begin{lstlisting}[language=Fortran]
&params
  ts%lat_file = "small_ring.bmad"
  ts%dat_out_file = "da_tune_scan.dat"
  
  ts%Q_a0 = 0.100
  ts%Q_a1 = 0.400
  ts%dQ_a = 0.005
  
  ts%Q_b0 = 0.490
  ts%Q_b1 = 0.700
  ts%dQ_b = 0.005
  
  ts%Q_z0 = 0.050
  ts%Q_z1 = 0.050
  ts%dQ_z = 0.01
  
  ts%n_turn = 1000
  ts%rf_on = T
  
  ts%min_angle = 0
  ts%max_angle = 3.1415926
  ts%n_angle = 37
  ts%x_init = 0.005
  ts%y_init = 0.005
  
  bmad_com%radiation_damping_on = F
  bmad_com%radiation_fluctuations_on = F
/
\end{lstlisting}

\subsection{Запуск}

\begin{verbatim}
$ cd bsim/da_tune_scan/example
$ da_tune_scan da_tune_scan.init
\end{verbatim}

\section{Выходные данные}

Формат выходного файла:
\begin{verbatim}
# ja jb jz Q_a Q_b Q_z DA_area DA_x_max DA_y_max n_angle_survived calc_successful
 0  0  0  0.100 0.490 0.050 1.23e-4 1.5e-2 1.2e-2 35 T
 1  0  0  0.105 0.490 0.050 1.18e-4 1.4e-2 1.1e-2 34 T
...
\end{verbatim}

\section{Замечания}

\begin{itemize}
    \item Время расчета пропорционально произведению количества точек сетки на количество угловых направлений
    \item Вблизи резонансов настройка частот может быть невозможна -- это нормальная ситуация
    \item Для ускорения расчета можно уменьшить \texttt{n\_angle} или увеличить шаг сетки
    \item Необходимо наличие апертур в lattice файле для корректного расчета DA
\end{itemize}

\end{document}
